\subsection{A noisy code is a genuine full-AA obstruction}
\label{threshold:noisy-code}

Let $m\ge2$ be even, $n=2^m$, and index coordinates by
$z\in\mathbb F_2^m$.  Put $c_a(z)=(-1)^{a\cdot z}$ and
\[
 h(a)=\sum_{j=1}^{m/2}a_{2j-1}a_{2j}\pmod 2,
 \qquad r=\Pr[h(A)=1]=\frac{1-n^{-1/2}}2.
\]
Draw $A,B$ uniformly and independently from $\mathbb F_2^m,\mathbb F_2$,
and independent signs $\eta_z$ with $\mathbb E\eta_z=\rho$, where
$0<\rho<1$.  The noisy affine-code law $\nu_{n,\rho}$ is
\[
 X_z=(-1)^B c_A(z)\eta_z.
\]
Define on the entire cube
\[
 P(x)=\sum_{a:h(a)=1}\left(\frac1n\sum_z x_zc_a(z)\right)^2.
\]
Parseval gives $0\le P\le1$.  A single phase query, followed by the
Walsh transform and the test $h(a)=1$, has acceptance probability $P(x)$
for every input.  This is the hidden-character quantum procedure of
\href{https://people.eecs.berkeley.edu/~vazirani/pubs/bv.pdf}
{Bernstein and Vazirani}, with an additional classical decision rule.

\begin{theorem}[Noise-robust one-query separation]
\label{threshold:noisy-query-theorem}
Every classical tree $D$ of depth at most $t$ satisfies
\[
 \boxed{\displaystyle
 \mathbb E_{\nu_{n,\rho}}|D-P|
 \ge \frac{\rho^2}{2}\bigl(1-2^{t-m/2}\bigr)
 -\sqrt{\frac{2r(1-r)(1-\rho^4)}{n}}.}
 \tag{NQ}
\]
This holds also with randomized queries or a randomized real output.
Consequently $\mathcal A(\nu_{\rho})$ fails for every fixed $\rho>0$:
one quantum query requires $\Omega(\log n)$ classical queries at a
fixed positive accuracy depending only on $\rho$.
\end{theorem}

\begin{proof}
The projection calculation of Section~\ref{proof:higher_order_obstructions}
gives
\[
 \mathbb E[P\mid A,B]=(1-\rho^2)r+\rho^2h(A)=:G(A),
 \qquad
 \Var(P)=r(1-r)\left[\rho^4+\frac{2(1-\rho^4)}n\right].
\]
The first identity can also be read directly from the quadratic form:
diagonal terms are unchanged, and distinct-coordinate terms acquire
the factor $\rho^2$.  Since $\Var(G)=\rho^4r(1-r)$, the law of total
variance and Cauchy--Schwarz imply
\[
 \mathbb E|P-G(A)|
 \le\sqrt{\mathbb E\Var(P\mid A,B)}
 =\sqrt{2r(1-r)(1-\rho^4)/n}.
 \tag{NC}
\]

Fix $B$, the entire noise vector $\eta$, and any private random seed of
the classical algorithm.  A query to $X_z$ can then be answered from
one query to the clean character $c_A(z)$, by multiplying the answer by
the fixed sign $(-1)^B\eta_z$.  Even granting these fixed signs to the
algorithm cannot invalidate a lower bound.  A depth-$t$ transcript
therefore conditions uniform $A$ on an affine subspace of codimension
at most $t$.

The Boolean sign $(-1)^{h(A)}$ has all Fourier coefficients of magnitude
$2^{-m/2}$.  Thus its conditional bias at every such leaf is at most
$2^{t-m/2}$, exactly as in
Section~\ref{qm:one-quantum-query-versus-logarithmically-many-classical-queries}.
The target $G(A)$ takes two values at distance $\rho^2$.  The minimum
conditional absolute loss of any real prediction is this distance
times the smaller conditional label probability.  Hence
\[
 \mathbb E|D-G(A)|\ge
 \rho^2(1-2^{t-m/2})/2.
\]
Averaging over the fixed signs and seeds, and then using (NC) and the
triangle inequality, proves (NQ).  For $t\le m/2-2$, the leading term
is at least $3\rho^2/8$; for sufficiently large $n$, the subtracted
term is at most $\rho^2/16$.  Accuracy $\rho^2/8$ therefore requires
$t>m/2-2$.
\end{proof}

The noise is sampled once as part of the input.  It is not resampled
at each oracle invocation.  This distinction is needed for the quantum
acceptance formula and for the classical simulation by a fixed noise
vector.  The result strengthens the earlier influence-only statement
for this same law to a query separation.
